\section{Introduction}
\label{sec:introduction}

Partisan gerrymandering through packing and cracking leaves a distinctive geometric
signature on the district vote-share distribution. Cracking spreads the target party's
voters across many districts at vote shares just below 50\%---winning voters concentrated
in losing positions, with all votes wasted. Packing concentrates the remaining voters in
a small number of landslide districts at 70--80\%---the excess votes above 50\%+1 wasted
as surplus. The combination produces a bimodal distribution: many districts clustered near
45\% Democratic (cracked) and a few districts clustered near 75\% Democratic (packed),
with the mean and median of Democratic-won and Democratic-lost districts pulled apart.

\citet{warrington2018} proposed Declination as a direct geometric measure of this pattern.
Rather than aggregating wasted votes (as EG does) or measuring the mean-median gap (as MM
does), Declination computes the angular geometry of the vote-share distribution: the angle
at which Republican-won districts tilt toward 50\% (cracking: opponents' wins cluster just
below 50\%) versus the angle at which Democratic-won districts tilt toward 50\% (packing:
opponents' wins cluster far above 50\%). The difference in these angles, normalized to the
range $[-1, 1]$, is the Declination.

This geometric approach has an appealing transparency: Declination can be read directly
from a plot of sorted district vote shares, without requiring computation of wasted votes,
median statistics, or swing models. A court or jury can see the asymmetry directly in the
graph---many districts clustered just below the 50\% line (cracked districts) versus few
districts far above it (packed districts)---and the Declination translates that visual
asymmetry into a scalar.

The legal status of Declination is, however, explicitly undeveloped. No federal or state
court has adopted it as evidence as of 2026. This paper documents that legal status
explicitly and explains the methodological critiques that may account for judicial
non-adoption. It also provides the first systematic analysis of Declination's behavior
under \textsc{bisect} algorithmic redistricting, establishing the neutral-map baseline
for comparison to enacted maps.

The paper's contribution is not to rehabilitate Declination for court use---its legal
status is what it is---but to provide a rigorous empirical foundation for expert witnesses
who may be asked about it in cross-examination, and to establish that \textsc{bisect}
maps produce near-zero Declination as a byproduct of compactness, providing context for
evaluating enacted maps' Declination values if courts eventually develop interest in the metric.

Section~\ref{sec:definition} defines Declination and documents the three methodological
critiques. Section~\ref{sec:behavior} explains the compactness mechanism. Section~\ref{sec:empirical}
presents empirical results. Section~\ref{sec:legal} documents the legal non-adoption record
through 2026. Section~\ref{sec:conclusion} concludes.
